% ID: SEL-D0058-OA-E02
% ESTADO: cotejada. Segunda pasada visual y borrador compilado revisados.
% Fuente: Selectividad/sel_2011_matematicas_aplicadas/221A-EXAMEN 5.pdf
% SHA-256: 1780d560801fb3001e146f12b3a37b29b566ca6576dc9cdafda07902d7dee982
% Página: 1; opción A; ejercicio original 2.

\begin{enumerate}[label=\alph*)]
\item \textbf{(1.25 puntos)} Halle el dominio, los puntos de corte con los ejes, y las asíntotas de la función
\[
f(x)=\frac{4x}{2x+1}.
\]
\item \textbf{(1.25 puntos)} Halle los intervalos de monotonía, los extremos relativos, los intervalos de curvatura y los puntos de inflexión de la función $g(x)=x^3+3x^2+3x$.
\end{enumerate}
